% !PRINTED 360
\weberpage{360}
Now the necessary and sufficient condition for a ternary quadratic form to split into two linear factors is that the determinant of the quadratic form vanish (Bd. I, \S.\,56).

If we express the forms $u, v$ in terms of $x_1, y_1$ and some third variable $y$, and denote the coefficients in $v$ and $u$ by $a_{ik}, b_{ik}$, so that $2a_{23}, 2b_{23}$ become the coefficients of $x_1 y_1$, we obtain the condition for the splitting of $v + 2\lambda u + \lambda^2 x_1 y_1$ in the form:
\weq{6}{\begin{vmatrix} a_{11}+2\lambda b_{11}, & a_{12}+2\lambda b_{12}, & a_{13}+2\lambda b_{13} \\ a_{21}+2\lambda b_{21}, & a_{22}+2\lambda b_{22}, & a_{23}+2\lambda b_{23}+\tfrac{1}{2}\lambda^2 \\ a_{31}+2\lambda b_{31}, & a_{32}+2\lambda b_{32}+\tfrac{1}{2}\lambda^2, & a_{33}+2\lambda b_{33} \end{vmatrix} = 0,}
and this is an equation of degree $5$ with respect to $\lambda$, which therefore gives us five decompositions:
\[
x_2 y_2,\ x_3 y_3,\ x_4 y_4,\ x_5 y_5,\ x_6 y_6.
\]
The coefficients in $f$ can, as (6) shows, be so disposed that all the functions $x_i y_i$ thus determined are different from one another, and from this it follows, as shown above, that they remain different so long as the curve $f$ has no singular points. Hence the following theorem holds:

1. \gesperrt{To every pair of double tangents $x_1 y_1$ there belong five further pairs $x_i y_i$ of such a kind that the eight points of contact of $x_1 y_1 x_i y_i$ lie on a conic section.}

The six pairs of double tangents that are determined in this way:
\[
x_1 y_1,\ x_2 y_2,\ x_3 y_3,\ x_4 y_4,\ x_5 y_5,\ x_6 y_6,
\]
we shall call a \gesperrt{Steiner complex}\footnote{Hitherto the expression \gesperrt{Steiner group} has commonly been used for this. But since the word ``group'' has an entirely definite other meaning in algebra, we prefer to avoid this expression here.}.

Let us consider the pairs $x_1 y_1, x_2 y_2, x_3 y_3$ of such a complex, and set
\[
f = x_1 y_1 x_2 y_2 - u_1^2 = x_1 y_1 x_3 y_3 - u_2^2,
\]
%CONT>>
